%% Springer Nature LaTeX template — sn-jnl (v3.1, December 2024)
%% Compile: pdflatex → bibtex → pdflatex → pdflatex

\documentclass[pdflatex,sn-mathphys-num]{sn-jnl}

%% ─── Packages ───────────────────────────────────────────────────────────────
\usepackage{graphicx}
\usepackage{amsmath}
\usepackage{amssymb}
\usepackage{booktabs}
\usepackage{multirow}
\usepackage{algorithm}
\usepackage{algpseudocode}
\usepackage{microtype}
\hypersetup{
    colorlinks=true,
    citecolor=black,
    linkcolor=black,
    urlcolor=black
}

%% ─── Title ──────────────────────────────────────────────────────────────────
\title[Context-Aware False Alarm Reduction for Violence Detection]{%
  Context-Aware False Alarm Reduction for Violence Detection
  in Surveillance Videos Using a Lightweight Gating Module}

%% ─── Authors ────────────────────────────────────────────────────────────────
\author*[1]{\fnm{Viet Hung} \sur{Ha}\orcidA{}}\email{hungha.060963@gmail.com}
\author[1]{\fnm{Viet Nhan} \sur{Nguyen}\orcidB{}}\email{nhannvse201082@fpt.edu.vn}
\author[1]{\fnm{Thai Kiet} \sur{Nguyen}\orcidC{}}\email{kietntse200734@fpt.edu.vn}
\author[1]{\fnm{Bao Nguyen} \sur{Tran}\orcidD{}}\email{nguyentbse201012@fpt.edu.vn}

\newcommand{\orcidA}{~\href{https://orcid.org/0009-0005-9924-4498}%
  {\textcolor{black}{[0009-0005-9924-4498]}}}
\newcommand{\orcidB}{~\href{https://orcid.org/0009-0007-7770-6449}%
  {\textcolor{black}{[0009-0007-7770-6449]}}}
\newcommand{\orcidC}{~\href{https://orcid.org/0009-0004-7609-9622}%
  {\textcolor{black}{[0009-0004-7609-9622]}}}
\newcommand{\orcidD}{~\href{https://orcid.org/0009-0006-9895-9827}%
  {\textcolor{black}{[0009-0006-9895-9827]}}}

\affil*[1]{%
  \orgname{FPT University},
  \orgaddress{%
    \city{Ho Chi Minh City},
    \country{Vietnam}}}


%% ─── Abstract ───────────────────────────────────────────────────────────────
\abstract{%
Violence detectors deployed on CCTV surveillance systems suffer from
unacceptably high false positive rates in contextually ambiguous
scenes---crowded gatherings, sporting events, or dimly lit
environments---where vigorous but benign motion closely resembles genuine
violent incidents. This problem, known as \emph{alert fatigue}, causes
security operators to ignore repeated erroneous alarms, ultimately
undermining the system's safety purpose. Most existing deep-learning
approaches pursue ever-higher classification accuracy by replacing the
visual backbone, yet they reason solely over appearance and motion cues and
seldom treat the false positive rate~(FPR) as an explicit design objective.
We take a complementary, deployment-driven view and ask: given any
reasonable detector, how can its false alarms be reduced cheaply using
freely available scene context?
To this end we propose a lightweight, model-agnostic two-stage framework
that augments a pre-trained violence detector with a \emph{Context Gating
Module}~(CGM), suppressing false alarms without retraining the backbone.
In the first stage, a frozen X3D-S backbone~\cite{fan2020x3d} fine-tuned
on RWF-2000~\cite{cheng2021rwf} provides an initial violence
probability~$p_{\text{base}}$. In the second stage, three parallel context
streams are extracted directly from raw video without any additional
annotation: crowd density~(YOLOv8n), lighting condition~(OpenCV), and a
scene-level motion-chaos descriptor capturing direction entropy and motion
synchrony~(Farneback optical flow). The resulting 13-dimensional context
vector drives the CGM---two compact MLPs totalling only 962 trainable
parameters. An MLP-gate produces an interpretable attention
weight~$\alpha\in[0,1]$ encoding how much the base detector should be
trusted for a given clip, while an MLP-ctx produces a context-corrected
probability~$p_{\text{ctx}}$, yielding the calibrated output
$p_{\text{final}}=\alpha\cdot p_{\text{base}}+(1-\alpha)\cdot p_{\text{ctx}}$.
On RWF-2000, the CGM reduces the false positive rate from $0.153$ to $0.113$
---a $26.1\%$ relative reduction---while raising accuracy from
$85.95\%$ to $87.63\%$ and keeping the false negative rate essentially
unchanged. A per-stream ablation identifies the lighting stream as the
strongest single context cue and shows that the full three-stream
combination outperforms any individual stream, indicating a synergistic
effect. We further examine whether the framework transfers across
distributions by evaluating on the RLVS dataset~\cite{soliman2019rlvs}.}

\keywords{Violence detection, False alarm reduction, Context-aware
  calibration, Gating module, X3D, Surveillance video, Optical flow,
  Crowd density, Cross-dataset generalization, Lightweight deep learning}

%% ─── Document ───────────────────────────────────────────────────────────────
\begin{document}

\maketitle

\section{Introduction}
\label{sec:intro}

Automated violence detection in CCTV surveillance is a critical component
of modern public-safety infrastructure. Recent deep-learning approaches
achieve strong accuracy on standard benchmarks, ranging from efficient
3D convolutional networks~\cite{fan2020x3d} and two-stream
designs~\cite{cheng2021rwf} to transformer- and state-space
architectures~\cite{senadeera2024cuenet,senadeera2025videomamba}.
Yet a persistent practical problem remains: excessive false positive
alerts in contextually ambiguous scenes such as crowded gatherings,
sporting events, or dimly lit environments. Repeated erroneous alarms
cause \emph{alert fatigue}~\cite{islam2023survey}, eroding operator trust
and ultimately undermining the system's safety purpose.

The root cause is twofold. First, the dominant research objective is
\emph{accuracy}, not the false positive rate~(FPR); a detector reporting
$88\%$ accuracy can still raise an alarm on a large fraction of benign
clips, which is the metric that actually governs deployment cost. Second,
existing detectors reason almost exclusively over appearance and motion
cues and therefore cannot separate vigorous-but-benign crowd activity from
genuine violence. What distinguishes these situations is environmental
\emph{context}---how crowded the scene is, how well it is lit, and whether
motion is spatially coherent (as in coordinated crowd movement or camera
shake) rather than chaotic.

Prior work leaves this gap open. Backbone-centric
methods~\cite{senadeera2024cuenet,senadeera2025videomamba,visafe2025}
improve accuracy and efficiency by redesigning the network, but neither
target FPR nor exploit scene context. Approaches that do model
context~\cite{fistfight2025} rely on expensive, annotation-dependent
pipelines such as pose estimation and tracking. Gating has been shown to be
an effective fusion primitive~\cite{senadeera2025videomamba}, but the gate
is embedded \emph{inside} the backbone and trained end to end, so it cannot
be reused with a different detector. A survey of over 200
papers~\cite{islam2023survey} confirms that no existing work explicitly
optimises FPR through lightweight, annotation-free contextual
post-processing applied on top of a frozen detector.

We address this gap with a lightweight, model-agnostic two-stage framework.
A frozen X3D-S backbone~\cite{fan2020x3d} fine-tuned on
RWF-2000~\cite{cheng2021rwf} provides an initial violence
probability~$p_{\text{base}}$. A \emph{Context Gating Module}~(CGM) of only
962 parameters then calibrates this probability using three automatically
extracted context streams---crowd density (YOLOv8n), lighting condition
(OpenCV), and a scene-level motion-chaos descriptor (Farneback optical
flow)---producing a final calibrated prediction
$p_{\text{final}}=\alpha\cdot p_{\text{base}}+(1-\alpha)\cdot p_{\text{ctx}}$
(Eq.~\ref{eq:gating}). Because the backbone is frozen and the gate is
external, the module attaches to any probabilistic detector at negligible
cost.

Our study is organised around three research questions:
\begin{itemize}
  \item \textbf{RQ1.} Does augmenting X3D-S with the Context Gating Module
        reduce the false positive rate compared with X3D-S alone, without
        increasing the false negative rate?
  \item \textbf{RQ2.} Which context stream---crowd, lighting, or
        motion-chaos---contributes most to the reduction in false alarms?
  \item \textbf{RQ3.} Does the framework generalise across datasets, i.e.\
        does a CGM trained on RWF-2000 still reduce false alarms when
        transferred to the RLVS dataset~\cite{soliman2019rlvs}?
\end{itemize}

\paragraph{Contributions.}
(i)~A \emph{model-agnostic}, two-stage framework that treats false alarm
reduction as an explicit objective and requires no retraining of the
underlying detector.
(ii)~An \emph{annotation-free} three-stream context extraction pipeline
(crowd, lighting, motion-chaos) that supplies information complementary to
the visual backbone.
(iii)~A compact \emph{Context Gating Module} (962 parameters) whose gating
weight~$\alpha$ yields an interpretable, per-clip measure of how much the
base detector is trusted.
(iv)~A systematic empirical study on RWF-2000---including a per-stream
ablation and a cross-dataset evaluation on RLVS---that quantifies the FPR
reduction, identifies lighting as the strongest single context cue, and
analyses the contribution of the motion-chaos stream.

The paper proceeds as follows: Section~\ref{sec:related} reviews related
work; Section~\ref{sec:method} details the methodology;
Sections~\ref{sec:experiments}--\ref{sec:results} present experiments and
results; Section~\ref{sec:conclusion} concludes.

%% ════════════════════════════════════════════════════════════════════════════
\section{Related Work}
\label{sec:related}

%% ════════════════════════════════════════════════════════════════════════════
\section{Methodology}
\label{sec:method}

%% ════════════════════════════════════════════════════════════════════════════
\section{Experiments}
\label{sec:experiments}

%% ════════════════════════════════════════════════════════════════════════════
\section{Results and Discussion}
\label{sec:results}

%% ════════════════════════════════════════════════════════════════════════════
\section{Conclusion}
\label{sec:conclusion}

%% ─── Bibliography ───────────────────────────────────────────────────────────
%% Note: bibliographystyle is set automatically by the sn-jnl class
%% (sn-mathphys-num); do not redeclare it here.
\bibliography{references}

\end{document}